% !PRINTED 591
\weberpage{591}
Since the congruence
\[
\frac{2\varrho}{\mu} + \frac{2\varrho_1}{\mu_1} \equiv 1\ (\mathrm{mod}\ 2,\,\omega)
\]
is not possible, as one shows exactly as above (p.\ 589), the following is
necessary and sufficient for equation (3):
\weq{4}{\frac{2(\mu_1\varrho - \mu\varrho_1)}{\mu\mu_1} \equiv 0\ (\mathrm{mod}\ 2,\,\omega).}
From this it follows:
\weq{5}{\mu\varrho_1 \equiv \mu_1\varrho\ (\mathrm{mod}\ \mu\mu_1),}
and this is possible only if
\weq{6}{\varrho \equiv 0\ (\mathrm{mod}\ \mathfrak{a}), \quad \varrho_1 \equiv 0\ (\mathrm{mod}\ \mathfrak{a}_1).}
If conversely the conditions (5) and (6) are satisfied, then:
\[
\frac{2\varrho}{\mu} - \frac{2\varrho_1}{\mu_1}
\]
is an integer and equation (3) is satisfied.

Now take an integer in $\Omega$ divisible by $\mathfrak{a}$:
\weq{7}{\alpha = \mathfrak{a}\mathfrak{c},}
in which $\mathfrak{c}$ is relatively prime to $\mu$ and $\mu_1$. If then $\varrho$ is
divisible by $\mathfrak{a}$, one can determine an integer $\xi$ with respect to the
modulus $\mathfrak{m}$ from the congruence
\[
\varrho \equiv \alpha\xi\ (\mathrm{mod}\ \mu)
\]
(Bd.\ II, \S\,166, 7.); then
\weq{8}{\alpha_1 = \frac{\alpha\mu_1}{\mu} = \mathfrak{a}_1\mathfrak{c}}
is an integer divisible by $\mathfrak{a}_1\mathfrak{c}$ and by no other ideal,
and from (5) there results
\[
\varrho_1 \equiv \alpha_1\xi\ (\mathrm{mod}\ \mu_1).
\]
5. \gesperrt{We thus obtain the common roots of
$\boldsymbol{A}_\mu(x), \boldsymbol{A}_{\mu_1}(x)$ in the form:}
\weq{9}{\boldsymbol{s}\!\left(\frac{2\alpha\xi}{\mu}\right) = \boldsymbol{s}\!\left(\frac{2\alpha_1\xi}{\mu_1}\right),}
\gesperrt{where $\xi$ runs through a complete system of residues with respect to the modulus $\mathfrak{m}$.}

\gesperrt{Here $\mathfrak{m}$ can denote any arbitrary odd ideal in $\Omega$.}

For to every such ideal one can choose two numbers $\mu, \mu_1$
of the form (2), and then choose $\omega$ among the equivalent
numbers in such a way that $\mathfrak{a}$ becomes odd and relatively prime to $\mu\mu_1$.
